\section{Kỹ thuật context}

\subsection{Câu hỏi nghiên cứu}

\textit{Làm thế nào biến Character State + Memory + Relationship + World + Scenario thành context tối ưu cho LLM?}

\subsection{So sánh Methods}

\begin{table}[H]
\centering
\caption{Các Phương pháp Context Engineering}
\label{tab:context-methods}
\begin{tabular}{lllll}
\toprule
\textbf{Method} & \textbf{Token Reduction} & \textbf{Accuracy} & \textbf{Latency} & \textbf{Use Case} \\
\midrule
Naive concatenation & 0\% & Baseline & Baseline & Trường hợp đơn giản \\
LongLLMLingua & 60\% & $-2\%$ & $+20\%$ & Token-constrained \\
Self-RAG & 40\% & $+5\%$ & $+50\%$ & Retrieval-needed \\
GraphRAG & 50\% & $+8\%$ & $+80\%$ & Complex reasoning \\
\bottomrule
\end{tabular}
\end{table}

\subsection{Chiến lược Compilation}

\textbf{Thứ tự ưu tiên}: Identity $>$ Personality $>$ Memory $>$ Relationship $>$ World $>$ Scenario

\textbf{Nén}: Kiểu LongLLMLingua cho components ưu tiên thấp
\textbf{Giữ nguyên}: Chi tiết đầy đủ cho Identity và Personality

% =============================================================================
% 14. MACHINE LEARNING \& DEEP LEARNING